\documentclass[hyphens,aspectratio=169,dvipsnames]{beamer}
\usepackage{graphicx}
\usepackage{xcolor}
\usepackage{amssymb}
\usepackage{amsmath}
\usepackage{mathtools}
\usepackage{textcomp}
\usepackage{moresize}
\usepackage{framed}
\usepackage{minted}
\usepackage{relsize}
\usepackage{emoji}
\usepackage{tikz}
\usepackage{comment}
\usetikzlibrary{shapes.geometric, arrows, positioning}
\usetheme{Berlin}
\usecolortheme[RGB={0,95,47}]{structure}
\beamertemplatenavigationsymbolsempty

\makeatletter
\AtBeginEnvironment{minted}{\dontdofcolorbox}
\def\dontdofcolorbox{\renewcommand\fcolorbox[4][]{##4}}
\makeatother

\newcommand{\textpf}[1]{\texttt{\color{black}\fcolorbox{lightgray}{lightgray}{#1}}}

\begin{document}

\begin{frame}[fragile]{Tell me what to type!}
    Let's try a C program that runs \textpf{ls}, then runs \textpf{ls /}.
    % The goal is for the students to end up with something like `ls_twice.c` (found in this directory).
\end{frame}

\begin{frame}[fragile]{Let's \textpf{strace} it}
    \begin{itemize}
        \pause \item Notice anything different? % The order of execution has changed, due to the overhead of ptrace.
    \end{itemize}
\end{frame}

\begin{frame}[fragile]{\textpf{waitpid}}
    \begin{itemize}
        \pause \item The \textpf{waitpid} system call allows processes to wait until events occur in other processes. It takes 3 arguments:
            \begin{itemize}
                \pause \item \textpf{pid\_t pid}, the \textpf{pid} of the process to wait on.
                \pause \item \textpf{int *wstatus}. This one is complicated. Wait for the next slide :)
                \pause \item \textpf{int options}, a set of flags. Just pass \textpf{0} for this argument.
            \end{itemize}
        \pause \item The returned value from \textpf{waitpid} is the PID of the \textpf{wait}ed-upon child, or \textpf{-1} on error.
    \end{itemize}
\end{frame}

\begin{frame}[fragile]{\textpf{int *wstatus}}
    \begin{itemize}
        \pause \item This is an ``out-argument." That is, it's a pointer argument that's used as an additional returned value.
        \pause \item Upon success, \textpf{waitpid} writes details about the \textpf{wait}ed-upon process into \textpf{*wstatus}.
        \pause \item You can access these details using function-like macros provided in \textpf{sys/wait.h}.
        \pause \item e.g., check \textpf{if (WIFEXITED(wstatus))} to test whether the \textpf{wait}ed-upon process has exited.
        \pause \item If it has exited, then you can obtain its exit status with \textpf{WEXITSTATUS(wstatus)}.
    \end{itemize}
\end{frame}

\begin{frame}[fragile]{Tell me what to type, again!}
    Let's fix our code from earlier.
    % The goal is for the students to end up with something like `ls_twice_fixed.c` (found in this directory).
\end{frame}

\begin{frame}[fragile]{The CWD (``current working directory")}
    \begin{itemize}
        \pause \item When you run \textpf{ls} in the shell, how does it know which directory to list?
        \pause \item Each process has a CWD, which is, by default, the directory from which that process was executed.
        \pause \item A process can change its CWD by issuing a \textpf{chdir} system call.
        \pause \item When you run \textpf{cd} in your shell, you're asking the shell to make a \textpf{chdir} syscall.
        \pause \item Demo! % Use `strace --trace=chdir bash` and then `cd` around
    \end{itemize}
\end{frame}

\begin{frame}[fragile]{\textpf{chdir}}
    \begin{itemize}
        \pause \item \textpf{chdir} is a system call that changes the CWD.
        \pause \item It takes a path string as an argument,
        \pause \item returns 0 on success,
        \pause \item and negative on error.
    \end{itemize}
\end{frame}

\begin{frame}[fragile]{Signals}
    \begin{itemize}
        \pause \item A signal is a little ``ping" that processes can send to each other.
        \pause \item When a process receives a signal, it will automatically run a special function associated with that signal, called a signal handler.
        \pause \item There are $\sim 32$ kinds of signals (listed in \textpf{man 7 signal}). e.g.,
            \begin{itemize}
                \pause \item \textpf{SIGSEGV}: The signal that the kernel sends to your process when you segfault.
                    \begin{itemize}
                        \pause \item Default handler: Terminate the process.
                    \end{itemize}
                \pause \item \textpf{SIGILL}: The signal that the kernel sends to your process when you attempt to execute an invalid instruction.
                    \begin{itemize}
                        \pause \item Default handler: Terminate the process.
                    \end{itemize}
                \pause \item \textpf{SIGINT}: The signal that your process receives when you press Ctrl-C in the terminal.
                    \begin{itemize}
                        \pause \item Default handler: Terminate the process.
                    \end{itemize}
            \end{itemize}
        \pause \item The different kinds of signals are defined as constants in \textpf{signal.h}.
    \end{itemize}
\end{frame}

\begin{frame}[fragile]{Sending Signals}
    \begin{itemize}
        \pause \item To send a signal to a process, use \textpf{kill(2)} in C, or \textpf{kill(1)} in the shell.
        \pause \item \textbf{IMPORTANT: \textpf{kill} sends a signal! It doesn't necessarily kill the process!}
        \pause \item To the \textpf{man} page!
    \end{itemize}
\end{frame}

\begin{frame}[fragile]{Signal Handlers}
    \begin{itemize}
        \pause \item For most types of signals, a process may register a \textit{signal handler}: a function that is invoked when the signal is delivered.
        \pause \item A signal handler should accept one \textpf{int} argument, and return \textpf{void}.
        \pause \item The argument will contain the numeric value of the received signal type (which is safe to ignore if your signal handler handles only one type of signal).
        \pause \item You can register signal handlers with \textpf{signal(3)}.
        \pause \item To the \textpf{man} page!
    \end{itemize}
\end{frame}

\begin{frame}[fragile]{Signal Example}
    Demo!
    % simple_signal.c
\end{frame}


\begin{frame}[fragile]{The Dangers of Signals!}
    Demo!
    % malloc_loop.c
\end{frame}

\begin{frame}[fragile]{Activity!}
    Write a C program that
    \begin{enumerate}
        \pause \item parses a hex number out of \textpf{argv[1]} using \textpf{strtoull(3)},
        \pause \item then writes a single \textpf{'a'} byte into that address,
        \pause \item and exits successfully only when that write causes a \textpf{SIGSEGV} to be delivered to your process.
    \end{enumerate}
\end{frame}

\end{document}
